% ==================== 八、实验总结与不足 ====================
\section{实验总结与不足}

\subsection{设计产物概览}

本报告各章给出的设计产物如下，便于按图查阅：

\begin{table}[H]
\centering
\caption{各章的设计产物}
\tabfont
\begin{tabularx}{\textwidth}{@{}L{2.2cm}L{4.2cm}L{2.6cm}X@{}}
\toprule
章节 & 设计产物 & 规模 & 内容 \\
\midrule
第 2 章 & 功能划分与需求追溯 & 10 个模块、36 条需求 & 需求逐条落模块，跨模块与不覆盖的单独标注 \\
第 3 章 & 架构设计与模块交互 & 包图 / 组件图 / 部署图 & 模块依赖、交互链路与八条设计决策的代价 \\
第 4 章 & 数据库设计 & 22 张表 & 字段、索引、唯一约束与库存条件更新 \\
第 5 章 & UML 建模 & 16 张图 & 结构 4 张、行为 8 张、交互 4 张，均为 PlantUML 源码加导出图 \\
第 6 章 & 界面设计与功能描述 & 19 屏 & 四端各屏的元素、状态及对应需求 \\
第 7 章 & AI 协作过程记录 & 设计阶段 7 组，另有报告改稿 & 事前讨论、生成结果、优缺点与本人校改 \\
第 7 章末 & 人机分工流程 & 一条五步流程 & 从确定内容到核对结果均保留人的判断 \\
\bottomrule
\end{tabularx}
\end{table}

\subsection{主要设计结果}

本次完成的主要设计内容为：

\textbf{库存。} 日库存按房型和晚次记录，同房型请求先赋序再处理；下单采用条件更新与事务回滚，避免并发超卖（REQ-011）。并发抢房流程见第 5.6 节。

\textbf{状态。} 分别绘制订单与房间状态图，列出订单转换条件，非法转换统一拒绝（REQ-005、REQ-008）。

\textbf{界面。} 绘制 19 屏并建立界面与需求的对应表。09 屏展示三步入住；11 屏展示线下保留额度不足时拒绝开单，并提示引导客人线上预订或由店长先调整额度（REQ-018、REQ-033）。

这次设计中，我对「预订成功」和「分配房号」的区别理解得更清楚了。下单时按房型和日期占库存，入住前才给订单安排具体房间；因此类图中需要单独的占房记录，前台界面也要能查看哪些房间在相应日期可分配。把这两个步骤拆开后，库存模型、类图和入住流程之间的关系才说得通。

报告和界面初稿由 AI 辅助生成，我负责对照实验一需求核对设计内容。19 屏界面完成后，曾出现跨屏余量、取消文案与业务规则不一致的问题；我依据统一示例数据和需求编号提出修改意见，再检查改后的图稿。这个过程让我发现，界面和 UML 图除了展示方案，也能用来检查前面讨论的规则有没有真正落到每个操作上。AI 的具体使用与修改记录见第 7 章。

\subsection{不足与后续工作}

\textbf{两处不足。} 一是初版下单设计只做到「不超卖」，REQ-011 后半句「按服务端收到提交的顺序分配最后一间房」是在核对需求原文时才补上的，说明读需求时对边界条件不够敏感。二是类图与数据表的对应关系是事后逐类核对出来的，初稿漏了片区、排班、价格区间等五个实体，补表时才发现 \texttt{user} 与 \texttt{store} 的 \texttt{region\_id} 没有父表。

\textbf{部署与并发规模。} 设计按单个片区、六家门店的规模展开，应用为单体，同房型请求队列放在进程内。换成多实例部署时，队列须改为外部组件（如 Redis 或数据库队列表），否则服务重启会丢失排队顺序。

\textbf{数据规模。} 示例数据覆盖一个片区与固定日期段；跨片区汇总、半年绩效归档与历史数据迁移只给出了表结构，没有用真实数据验证过分区与索引是否够用。

\textbf{界面与接口。} 19 屏覆盖主流程与常见失败提示；支付失败后的重试、库存冲突、网络中断等页面没有逐一成屏，目前以提示文案和规则说明表达。模块的方法签名、返回结构与错误码分级也留到实现阶段再定。

\textbf{待确认需求。} REQ-035（排房时限）与 REQ-036（注册与登录）在实验一尚未确认，设计没有给出具体流程；确认后需要补相应的模块职责、界面和校验规则。
